% ==============================================================================
% Section 6: Swiss Multinational Enterprises Abroad
% ==============================================================================

\section{Swiss Multinational Enterprises Abroad}

The macroeconomic footprint of Swiss multinational enterprises (MNEs) extends far beyond direct financial capital stocks. The operational reality of Swiss globalization is captured by structural enterprise indicators: the expansion of foreign operating affiliates, aggregate global turnover, and the magnitude of personnel employed across international subsidiaries.

\subsection{Evolution of Foreign Subsidiaries (2015--2024)}

Despite periods of macroeconomic fragmentation and rising protectionism globally, Swiss multinational corporations have steadily expanded their network of foreign direct affiliates, as documented in Table~\ref{tab:subsidiaries_trend}.

\begin{table}[H]
\centering
\small
\caption{Trajectory of Non-Resident Subsidiaries of Swiss Multinational Enterprises (2015--2024)}
\label{tab:subsidiaries_trend}
\renewcommand{\arraystretch}{1.15}
\begin{tabularx}{0.85\textwidth}{c r r >{\raggedright\arraybackslash}X}
\toprule
\rowcolor{tableheadbg}
\textbf{Year} & \textbf{Number of Foreign Affiliates} & \textbf{Annual Growth (\%)} & \textbf{Strategic MNE Phase} \\
\midrule
2015 & 17,923 & --- & Baseline Network \\
\rowcolor{tablerowbg}
2016 & 17,982 & +0.3\% & Targeted Additions \\
2017 & 18,696 & +4.0\% & European Supply Chain Integration \\
\rowcolor{tablerowbg}
2018 & 19,295 & +3.2\% & North American Tech Acquisitions \\
2019 & 19,597 & +1.6\% & Pre-Pandemic Expansion High \\
\rowcolor{tablerowbg}
2020 & 19,571 & -0.1\% & Pandemic Operational Freeze \\
2021 & 20,087 & +2.6\% & Post-Lockdown Restructuring \\
\rowcolor{tablerowbg}
2022 & 21,045 & +4.8\% & Nearshoring \& Global Diversification \\
2023 & 21,251 & +1.0\% & Network Consolidation \\
\rowcolor{tablerowbg}
2024 & 21,265 & +0.1\% & Mature Global Footprint \\
\bottomrule
\end{tabularx}
\vspace{0.3em}
\footnotesize
\textit{Source: Swiss National Bank (SNB). Counts include majority-owned non-resident subsidiaries.}
\normalsize
\end{table}

Between 2015 and 2024, the total number of foreign operating subsidiaries maintained by Swiss MNEs expanded from $17,923$ to $21,265$, reflecting a net addition of \textbf{$3,342$ overseas corporate entities} (\textbf{$+18.7\pct$ cumulative growth}). This steady progression confirms that Swiss corporate leaders pursue long-term structural internationalization rather than short-term opportunistic cross-border ventures.

\subsection{Employment in Swiss Subsidiaries Abroad}

The most compelling manifestation of Switzerland's global industrial reach is the scale of human capital employed outside Switzerland by Swiss parent enterprises (Table~\ref{tab:employment_trend}).

\begin{table}[H]
\centering
\small
\caption{Global Workforce Employed in Swiss Foreign Subsidiaries Abroad (2015--2024)}
\label{tab:employment_trend}
\renewcommand{\arraystretch}{1.15}
\begin{tabularx}{0.85\textwidth}{c r r >{\raggedright\arraybackslash}X}
\toprule
\rowcolor{tableheadbg}
\textbf{Year} & \textbf{Headcount (in Thousands)} & \textbf{Annual Growth (\%)} & \textbf{Macroeconomic Dynamic} \\
\midrule
2015 & 2,033.6 & --- & Baseline Global Labor Pool \\
\rowcolor{tablerowbg}
2016 & 2,066.5 & +1.6\% & Organic Overseas Expansion \\
2017 & 2,101.6 & +1.7\% & Industrial Automation Growth \\
\rowcolor{tablerowbg}
2018 & 2,167.2 & +3.1\% & Broad International Hiring \\
2019 & 2,130.8 & -1.7\% & Structural Productivity Optimization \\
\rowcolor{tablerowbg}
2020 & 2,099.1 & -1.5\% & Pandemic Operational Layoffs \\
2021 & 2,238.4 & +6.6\% & Rapid Rebound in Overseas Operations \\
\rowcolor{tablerowbg}
2022 & 2,358.0 & +5.3\% & Strategic Nearshoring Hiring \\
2023 & 2,474.9 & +5.0\% & Global Service \& R\&D Expansion \\
\rowcolor{tablerowbg}
2024 & 2,483.5 & +0.3\% & All-Time High Employment Level \\
\bottomrule
\end{tabularx}
\vspace{0.3em}
\footnotesize
\textit{Source: Swiss National Bank (SNB). Staff employed in non-resident enterprises controlled by Swiss MNEs.}
\normalsize
\end{table}

\begin{remarkbox}[The $4.4\times$ Global Labor Multiplier of Swiss Enterprise]
In 2024, non-resident affiliates of Swiss MNEs employed \textbf{$2,483,500$ people} worldwide. To put this figure into perspective:
\begin{itemize}[leftmargin=1.2em, itemsep=2pt]
    \item The overseas workforce of Swiss multinationals is nearly \textbf{three times larger} than Switzerland's entire domestic manufacturing employment base ($\approx 650,000$ persons).
    \item Compared to the $560,964$ domestic personnel employed by Swiss parent enterprises within Switzerland, Swiss MNEs sustain an \textbf{overseas-to-domestic employment multiplier of $4.43\times$}. For every single employee working at a multinational corporate headquarters or laboratory in Switzerland, Swiss firms support more than four employees across global manufacturing, distribution, and clinical facilities abroad.
\end{itemize}
\end{remarkbox}

% ------------------------------------------------------------------------------
% Pgfplots Figure 6.1: Dual-Axis Evolution of Subsidiaries and Employment
% Greyscale Lines with Distinct Markers and Y-Axes
% ------------------------------------------------------------------------------
\begin{figure}[H]
\centering
\begin{tikzpicture}
\begin{axis}[
    width=13.8cm,
    height=7.0cm,
    grid=both,
    grid style={line width=0.2pt, color=midgray!30},
    xlabel={\small\textbf{Year}},
    ylabel={\small\textbf{Foreign Affiliates (Count)}},
    xmin=2014.5, xmax=2024.5,
    ymin=17000, ymax=22000,
    xtick={2015, 2016, 2017, 2018, 2019, 2020, 2021, 2022, 2023, 2024},
    xticklabel style={font=\footnotesize},
    yticklabel style={font=\footnotesize},
    legend style={
        at={(0.03,0.95)},
        anchor=north west,
        font=\footnotesize,
        fill=white,
        draw=pureblack,
        line width=0.6pt
    },
    line width=1.3pt
]

    % Subsidiaries Curve (Pure Black Solid Line with Square Markers)
    \addplot[
        color=pureblack,
        mark=square*,
        mark size=2.2pt,
        line width=1.4pt
    ] coordinates {
        (2015, 17923)
        (2016, 17982)
        (2017, 18696)
        (2018, 19295)
        (2019, 19597)
        (2020, 19571)
        (2021, 20087)
        (2022, 21045)
        (2023, 21251)
        (2024, 21265)
    };
    \addlegendentry{Foreign Operating Subsidiaries (Count, Left Axis)}

\end{axis}

% Secondary Axis for Foreign Employment
\begin{axis}[
    width=13.8cm,
    height=7.0cm,
    xmin=2014.5, xmax=2024.5,
    ymin=1950, ymax=2600,
    axis y line*=right,
    axis x line=none,
    ylabel={\small\textbf{Global Workforce Abroad (Thousands of Persons)}},
    yticklabel style={font=\footnotesize},
    legend style={
        at={(0.03,0.82)},
        anchor=north west,
        font=\footnotesize,
        fill=white,
        draw=pureblack,
        line width=0.6pt
    },
    line width=1.3pt
]

    % Workforce Curve (Dark Gray Dashed Line with Filled Triangle Markers)
    \addplot[
        color=darkgraytext,
        dashed,
        mark=triangle*,
        mark size=2.5pt,
        line width=1.4pt
    ] coordinates {
        (2015, 2033.6)
        (2016, 2066.5)
        (2017, 2101.6)
        (2018, 2167.2)
        (2019, 2130.8)
        (2020, 2099.1)
        (2021, 2238.4)
        (2022, 2358.0)
        (2023, 2474.9)
        (2024, 2483.5)
    };
    \addlegendentry{Foreign Affiliate Employment (Thousands, Right Axis)}

\end{axis}
\end{tikzpicture}
\caption{Synchronous Expansion of Swiss MNE Foreign Operating Subsidiaries and Global Affiliate Employment (2015--2024). Source: SNB. Both indicators display robust resilience and structural expansion over the decade.}
\label{fig:subsidiary_expansion}
\end{figure}

\subsection{Geographic Distribution of Foreign Affiliates}

The regional dispersion of the $21,265$ foreign subsidiaries operating in 2024 illustrates Swiss multinationals' multi-hub global configuration:

\begin{itemize}[leftmargin=1.5em, itemsep=2.5pt]
    \item \textbf{Europe:} $12,892$ subsidiaries (\textbf{$60.6\pct$}), representing the primary theater of operation due to geographic contiguity, bilateral regulatory harmonization, and integrated pan-European supply networks.
    \item \textbf{Asia:} $3,184$ subsidiaries (\textbf{$15.0\pct$}), serving as the fastest-growing regional affiliate network, led by operational hubs in China, India, Singapore, and Japan.
    \item \textbf{North America:} $2,443$ subsidiaries (\textbf{$11.5\pct$}), characterized by high capital and R\&D intensity per affiliate in the United States and Canada.
    \item \textbf{Central and South America:} $1,497$ subsidiaries (\textbf{$7.0\pct$}), focused on agricultural processing, consumer markets, and manufacturing (Brazil, Mexico).
    \item \textbf{Africa and Oceania:} $834$ subsidiaries in Africa ($3.9\pct$) and $415$ in Oceania ($2.0\pct$), completing a comprehensive worldwide operational presence.
\end{itemize}

\begin{mnecase}{Nestlé's Decentralized International Value Chain Architecture}{nestle_architecture}
Headquartered in Vevey on Lake Geneva, \textbf{Nestlé S.A.} represents the archetypal Swiss multinational:
\begin{itemize}[leftmargin=1.2em, itemsep=2pt]
    \item \textbf{Global Production Proximity:} Operating over 340 factories across 77 countries, Nestlé produces locally to preserve freshness, adapt to regional tastes, and avoid cross-border perishable tariffs.
    \item \textbf{Centralized R\&D and IP Command:} Fundamental scientific research, nutritional biotechnology, and strategic patent ownership remain concentrated at the Nestlé Research Center in Lausanne, with proprietary knowledge licensed downstream to operating affiliates worldwide.
\end{itemize}
\end{mnecase}
